\documentclass[11pt,a4paper]{article}

% ---- Packages ----
\usepackage[utf8]{inputenc}
\usepackage[T1]{fontenc}
\usepackage{amsmath,amssymb,amsthm}
\usepackage{mathtools}
\usepackage[margin=1in,headheight=14pt]{geometry}
\usepackage{hyperref}
\usepackage{enumitem}
\usepackage{xcolor}
\usepackage{tcolorbox}
\usepackage{fancyhdr}
\usepackage{booktabs}

% ---- Hyperref ----
\hypersetup{
    colorlinks=true,
    linkcolor=red,
    citecolor=blue,
    urlcolor=blue
}

% ---- Header ----
\pagestyle{fancy}
\fancyhf{}
\fancyhead[L]{\textit{Lesson 1: Measure Spaces}}
\fancyhead[R]{\thepage}
\renewcommand{\headrulewidth}{0.4pt}

% ---- Theorem Environments ----
\theoremstyle{definition}
\newtheorem{definition}{Definition}[section]
\newtheorem{example}{Example}[section]
\newtheorem{exercise}{Exercise}[section]

\theoremstyle{plain}
\newtheorem{theorem}{Theorem}[section]
\newtheorem{lemma}[theorem]{Lemma}
\newtheorem{proposition}[theorem]{Proposition}
\newtheorem{corollary}[theorem]{Corollary}

\theoremstyle{remark}
\newtheorem{remark}{Remark}[section]

% ---- Colored Boxes ----
\tcbuselibrary{theorems,skins,breakable}

\newtcolorbox{keyresult}[1][]{
    colback=green!5!white,
    colframe=green!50!black,
    fonttitle=\bfseries,
    title=#1,
    breakable
}

\newtcolorbox{rlconnection}[1][]{
    colback=blue!5!white,
    colframe=blue!50!black,
    fonttitle=\bfseries,
    title=#1,
    breakable
}

\newtcolorbox{warning}[1][]{
    colback=red!5!white,
    colframe=red!50!black,
    fonttitle=\bfseries,
    title=#1,
    breakable
}

% ---- Operators ----
\newcommand{\R}{\mathbb{R}}
\newcommand{\N}{\mathbb{N}}
\newcommand{\Q}{\mathbb{Q}}
\newcommand{\Z}{\mathbb{Z}}
\newcommand{\Prob}{\mathbb{P}}
\newcommand{\E}{\mathbb{E}}
\newcommand{\indicator}{\mathbf{1}}
\newcommand{\powerset}{\mathcal{P}}
\newcommand{\Borel}{\mathcal{B}}

% ---- Title ----
\title{\textbf{Lesson 1: Measure Spaces}\\[10pt]
\large $\sigma$-Algebras, Measures, and Measurable Functions\\[5pt]
\normalsize Session 1 of 102 $\mid$ Phase I: Mathematical Foundations $\mid$ 8\,h\\
\normalsize Primary text: Durrett, \emph{Probability: Theory and Examples}, Chapter 1}
\author{Curriculum: A Rigorous Learning Path to Multi-Agent Deep RL}
\date{February 9, 2026}

\begin{document}

\maketitle

\begin{abstract}
This lesson develops the measure-theoretic foundations needed for rigorous probability
and, ultimately, for the convergence proofs of reinforcement learning algorithms.
We construct the framework of $\sigma$-algebras and measures from first principles,
prove Carath\'{e}odory's extension theorem, define measurable functions and their
approximation by simple functions, and connect everything to probability theory.
Every result is proved. The reader who masters this material will have the language
and tools to define conditional expectation, filtrations, and martingales rigorously
--- the machinery that drives all stochastic approximation proofs in RL.
\end{abstract}

\tableofcontents
\newpage

%% ============================================================
%% SECTION 1: INTRODUCTION
%% ============================================================
\section{Why Measure Theory?}
\label{sec:intro}

The convergence proofs of TD-learning, Q-learning, and policy gradient methods
rest on stochastic approximation theory, which in turn rests on martingale
convergence theorems. To state these theorems precisely, we need:

\begin{enumerate}[nosep]
    \item \textbf{$\sigma$-algebras} to model the information available at each time step (filtrations).
    \item \textbf{Measures} to assign probabilities to events in a mathematically consistent way.
    \item \textbf{Measurable functions} to formalize random variables and value functions.
    \item \textbf{Integration} (next lesson) to define expectations rigorously.
\end{enumerate}

Without measure theory, ``$\E[X \mid \mathcal{F}_t]$'' is just notation. With it,
conditional expectation becomes a well-defined random variable characterized by two
precise properties, and the entire convergence machinery follows.

\begin{rlconnection}[Why This Matters for RL]
The TD(0) convergence proof (Theorem 4.2 in the TD article) uses:
\begin{itemize}[nosep]
    \item The filtration $\mathcal{F}_t = \sigma(S_0, A_0, R_1, \ldots, S_t)$ --- a $\sigma$-algebra generated by random variables.
    \item The noise term $w_k(s)$ satisfying $\E[w_k(s) \mid \mathcal{F}_k] = 0$ --- conditional expectation with respect to a $\sigma$-algebra.
    \item The Robbins--Siegmund theorem, which requires a \emph{non-negative supermartingale} --- defined via $\sigma$-algebras and conditional expectation.
\end{itemize}
This lesson provides the language for all three.
\end{rlconnection}

%% ============================================================
%% SECTION 2: SIGMA-ALGEBRAS
%% ============================================================
\section{$\sigma$-Algebras}
\label{sec:sigma_algebras}

\subsection{Definition and First Properties}
\label{subsec:sigma_def}

Throughout, $\Omega$ denotes a non-empty set (the \emph{sample space}).

\begin{definition}[$\sigma$-Algebra]
\label{def:sigma_algebra}
A collection $\mathcal{F} \subseteq \powerset(\Omega)$ of subsets of $\Omega$ is a
\emph{$\sigma$-algebra} (or \emph{$\sigma$-field}) on $\Omega$ if:
\begin{enumerate}[nosep,label=(\roman*)]
    \item $\Omega \in \mathcal{F}$.
    \item If $A \in \mathcal{F}$, then $A^c := \Omega \setminus A \in \mathcal{F}$.
    \hfill(closure under complement)
    \item If $A_1, A_2, \ldots \in \mathcal{F}$, then $\bigcup_{n=1}^\infty A_n \in \mathcal{F}$.
    \hfill(closure under countable union)
\end{enumerate}
The pair $(\Omega, \mathcal{F})$ is called a \emph{measurable space}. Elements of
$\mathcal{F}$ are called \emph{measurable sets} (or \emph{events} in probability).
\end{definition}

\begin{remark}[Why Countable, Not Arbitrary?]
\label{rem:countable}
Axiom (iii) requires closure under \emph{countable} unions, not arbitrary ones.
Closure under arbitrary unions would force $\mathcal{F} = \powerset(\Omega)$ in
most interesting cases, which is too large: as we will see
(Theorem~\ref{thm:vitali}), the power set of $\R$ admits no translation-invariant
measure extending the notion of ``length.'' The restriction to countable operations
is precisely what makes non-trivial measures possible.
\end{remark}

\begin{proposition}[Basic Properties]
\label{prop:sigma_basic}
Let $\mathcal{F}$ be a $\sigma$-algebra on $\Omega$. Then:
\begin{enumerate}[nosep,label=(\alph*)]
    \item $\emptyset \in \mathcal{F}$.
    \item If $A_1, A_2, \ldots \in \mathcal{F}$, then $\bigcap_{n=1}^\infty A_n \in \mathcal{F}$.
    \hfill(closure under countable intersection)
    \item If $A, B \in \mathcal{F}$, then $A \setminus B \in \mathcal{F}$ and $A \triangle B \in \mathcal{F}$.
    \hfill(closure under difference and symmetric difference)
    \item $\mathcal{F}$ is closed under finite unions and finite intersections.
\end{enumerate}
\end{proposition}

\begin{proof}
\textbf{(a)} $\emptyset = \Omega^c \in \mathcal{F}$ by (i) and (ii).

\textbf{(b)} By De Morgan's law: $\bigcap_{n=1}^\infty A_n = \left(\bigcup_{n=1}^\infty A_n^c\right)^c$.
Each $A_n^c \in \mathcal{F}$ by (ii), so $\bigcup A_n^c \in \mathcal{F}$ by (iii),
and $\left(\bigcup A_n^c\right)^c \in \mathcal{F}$ by (ii).

\textbf{(c)} $A \setminus B = A \cap B^c$. Since $B^c \in \mathcal{F}$, and $A \cap B^c = (A^c \cup B)^c \in \mathcal{F}$ by (ii) and (iii). For symmetric difference: $A \triangle B = (A \setminus B) \cup (B \setminus A)$.

\textbf{(d)} For finite unions $A_1 \cup \cdots \cup A_n$, set $A_k = \emptyset$ for $k > n$
in (iii). Finite intersections follow from (b) similarly.
\end{proof}

\begin{example}[Extreme $\sigma$-Algebras]
\label{ex:extreme}
For any $\Omega$:
\begin{itemize}[nosep]
    \item The \emph{trivial $\sigma$-algebra} $\{\emptyset, \Omega\}$ is the smallest $\sigma$-algebra on $\Omega$.
    \item The \emph{power set} $\powerset(\Omega) = \{A : A \subseteq \Omega\}$ is the largest.
\end{itemize}
Every $\sigma$-algebra on $\Omega$ lies between these two extremes:
$\{\emptyset, \Omega\} \subseteq \mathcal{F} \subseteq \powerset(\Omega)$.
\end{example}

\begin{example}[Partition-Generated $\sigma$-Algebra]
\label{ex:partition}
Let $\Omega = \{1,2,3,4\}$ and consider the partition $\{\{1,2\}, \{3,4\}\}$.
The $\sigma$-algebra generated by this partition is:
\begin{equation}
    \mathcal{F} = \bigl\{\emptyset, \{1,2\}, \{3,4\}, \{1,2,3,4\}\bigr\}.
\end{equation}
This has 4 elements. In general, a partition of $\Omega$ into $k$ blocks generates
a $\sigma$-algebra with $2^k$ elements (all possible unions of blocks).
\end{example}

\begin{rlconnection}[Connection to RL]
In a finite MDP with state space $\mathcal{S} = \{s_1, \ldots, s_n\}$, the
``information contained in knowing $S_t$'' is the $\sigma$-algebra generated by $S_t$:
$\sigma(S_t) = \sigma\bigl(\{S_t = s_1\}, \ldots, \{S_t = s_n\}\bigr)$.
This is a partition-generated $\sigma$-algebra with $2^n$ events. Each event
$A \in \sigma(S_t)$ is a union of events $\{S_t = s_i\}$, so knowing which
event occurred tells you exactly which state you are in.
\end{rlconnection}

\subsection{Generated $\sigma$-Algebras}
\label{subsec:generated}

Given an arbitrary collection of subsets, we want the ``smallest'' $\sigma$-algebra
containing them. This is essential because we often specify a few key events
(or random variables) and need to derive the full $\sigma$-algebra they generate.

\begin{lemma}[Intersection of $\sigma$-Algebras]
\label{lem:intersection}
If $\{\mathcal{F}_\alpha\}_{\alpha \in I}$ is any (possibly uncountable) collection
of $\sigma$-algebras on $\Omega$, then $\bigcap_{\alpha \in I} \mathcal{F}_\alpha$
is also a $\sigma$-algebra on $\Omega$.
\end{lemma}

\begin{proof}
We verify the three axioms.
\begin{enumerate}[nosep,label=(\roman*)]
    \item $\Omega \in \mathcal{F}_\alpha$ for all $\alpha$, so $\Omega \in \bigcap \mathcal{F}_\alpha$.
    \item If $A \in \bigcap \mathcal{F}_\alpha$, then $A \in \mathcal{F}_\alpha$ for all $\alpha$, so $A^c \in \mathcal{F}_\alpha$ for all $\alpha$, so $A^c \in \bigcap \mathcal{F}_\alpha$.
    \item If $A_1, A_2, \ldots \in \bigcap \mathcal{F}_\alpha$, then for each $\alpha$, all $A_n \in \mathcal{F}_\alpha$, so $\bigcup A_n \in \mathcal{F}_\alpha$. Hence $\bigcup A_n \in \bigcap \mathcal{F}_\alpha$. \qedhere
\end{enumerate}
\end{proof}

\begin{definition}[Generated $\sigma$-Algebra]
\label{def:generated}
For any collection $\mathcal{C} \subseteq \powerset(\Omega)$, the $\sigma$-algebra
\emph{generated by} $\mathcal{C}$ is:
\begin{equation}
    \sigma(\mathcal{C}) := \bigcap\bigl\{\mathcal{G} : \mathcal{G} \text{ is a $\sigma$-algebra on } \Omega,\; \mathcal{C} \subseteq \mathcal{G}\bigr\}.
    \label{eq:generated}
\end{equation}
\end{definition}

\begin{remark}
This is well defined by Lemma~\ref{lem:intersection}: the intersection of
$\sigma$-algebras is a $\sigma$-algebra, and the collection is non-empty (it
includes $\powerset(\Omega)$). By construction, $\sigma(\mathcal{C})$ is the
\emph{smallest} $\sigma$-algebra containing $\mathcal{C}$: it contains $\mathcal{C}$,
and any $\sigma$-algebra containing $\mathcal{C}$ must contain $\sigma(\mathcal{C})$.
\end{remark}

\begin{definition}[Borel $\sigma$-Algebra]
\label{def:borel}
The \emph{Borel $\sigma$-algebra} on $\R$, denoted $\Borel(\R)$, is the $\sigma$-algebra
generated by the open sets of $\R$:
\begin{equation}
    \Borel(\R) := \sigma\bigl(\{U \subseteq \R : U \text{ is open}\}\bigr).
    \label{eq:borel}
\end{equation}
Elements of $\Borel(\R)$ are called \emph{Borel sets}.
\end{definition}

\begin{proposition}[Equivalent Generators of $\Borel(\R)$]
\label{prop:borel_generators}
The following collections all generate the same $\sigma$-algebra $\Borel(\R)$:
\begin{enumerate}[nosep,label=(\alph*)]
    \item Open sets.
    \item Closed sets.
    \item Open intervals $(a, b)$ with $a, b \in \Q$.
    \item Half-lines $(-\infty, x]$ with $x \in \R$.
    \item Half-lines $(-\infty, x)$ with $x \in \Q$.
\end{enumerate}
\end{proposition}

\begin{proof}
We show that each collection generates a $\sigma$-algebra containing all the others.

\textbf{(a) $\Leftrightarrow$ (b):} Every closed set is the complement of an open set
and vice versa. Since $\sigma$-algebras are closed under complement,
$\sigma(\text{open sets}) = \sigma(\text{closed sets})$.

\textbf{(c) $\Rightarrow$ (a):} Every open set $U \subseteq \R$ is a countable
union of open intervals with rational endpoints (since $\Q$ is dense: for each
$x \in U$, find rational $a_x < x < b_x$ with $(a_x, b_x) \subseteq U$, and
$U = \bigcup_{x \in U} (a_{x}, b_{x})$; by density and separability, this is a countable union). So $\sigma(\text{rational open intervals}) \supseteq \text{open sets}$.

\textbf{(a) $\Rightarrow$ (c):} Trivially, since every open interval with rational
endpoints is open.

\textbf{(d) $\Rightarrow$ (a):} Note $(a, b) = (-\infty, b) \setminus (-\infty, a]$
and $(-\infty, b) = \bigcup_{n=1}^\infty (-\infty, b - 1/n]$. So open intervals
are in $\sigma(\{(-\infty, x] : x \in \R\})$, and by (c) $\Rightarrow$ (a), we
get all open sets.

\textbf{(a) $\Rightarrow$ (d):} $(-\infty, x] = \R \setminus (x, \infty)$, and
$(x, \infty)$ is open, so $(-\infty, x] \in \Borel(\R)$.

\textbf{(d) $\Leftrightarrow$ (e):} $(-\infty, x] = \bigcap_{n=1}^\infty (-\infty, x + 1/n)$ (using rationals $x + 1/n$ for a sequence approaching from the right, taking $q_n \in \Q$ with $x < q_n < x + 1/n$). Conversely,
$(-\infty, x) = \bigcup_{n=1}^\infty (-\infty, x - 1/n]$.
\end{proof}

\begin{remark}
Proposition~\ref{prop:borel_generators}(d) is particularly important. The
characterization via half-lines $(-\infty, x]$ means that to check whether a
function $f: \Omega \to \R$ is measurable, it suffices to verify
$f^{-1}((-\infty, x]) = \{\omega : f(\omega) \leq x\} \in \mathcal{F}$ for all
$x \in \R$. This is the standard criterion used for random variables.
\end{remark}

\subsection{The $\pi$--$\lambda$ Theorem}
\label{subsec:pi_lambda}

The $\pi$--$\lambda$ theorem (Dynkin's theorem) is one of the most useful tools
in measure theory. It lets us prove that two measures agree on an entire
$\sigma$-algebra by checking agreement on a much smaller collection.

\begin{definition}[$\pi$-System]
\label{def:pi_system}
A collection $\mathcal{P} \subseteq \powerset(\Omega)$ is a \emph{$\pi$-system}
if it is closed under finite intersection: $A, B \in \mathcal{P}$ implies
$A \cap B \in \mathcal{P}$.
\end{definition}

\begin{definition}[$\lambda$-System (Dynkin System)]
\label{def:lambda_system}
A collection $\mathcal{L} \subseteq \powerset(\Omega)$ is a \emph{$\lambda$-system}
if:
\begin{enumerate}[nosep,label=(\roman*)]
    \item $\Omega \in \mathcal{L}$.
    \item If $A, B \in \mathcal{L}$ and $A \subseteq B$, then $B \setminus A \in \mathcal{L}$.
    \hfill(closure under proper difference)
    \item If $A_1 \subseteq A_2 \subseteq \cdots$ are in $\mathcal{L}$, then
    $\bigcup_{n=1}^\infty A_n \in \mathcal{L}$.
    \hfill(closure under increasing union)
\end{enumerate}
\end{definition}

\begin{remark}
A $\sigma$-algebra is both a $\pi$-system and a $\lambda$-system. The power of the
$\pi$--$\lambda$ theorem is that the converse also holds: a collection that is
both a $\pi$-system \emph{and} a $\lambda$-system is a $\sigma$-algebra.
\end{remark}

\begin{lemma}[$\pi + \lambda = \sigma$]
\label{lem:pi_lambda_sigma}
If $\mathcal{F}$ is both a $\pi$-system and a $\lambda$-system, then $\mathcal{F}$ is a $\sigma$-algebra.
\end{lemma}

\begin{proof}
We verify the three $\sigma$-algebra axioms.

\textit{(i)} $\Omega \in \mathcal{F}$ by the $\lambda$-system property.

\textit{(ii)} If $A \in \mathcal{F}$, then $A \subseteq \Omega$ and both are in $\mathcal{F}$, so $A^c = \Omega \setminus A \in \mathcal{F}$ by the proper difference property.

\textit{(iii)} Let $A_1, A_2, \ldots \in \mathcal{F}$. Define $B_n = \bigcup_{k=1}^n A_k$.
We need to show $B_n \in \mathcal{F}$ for each $n$ (then $\bigcup A_n = \bigcup B_n$
is an increasing union in $\mathcal{F}$).

For $n = 1$: $B_1 = A_1 \in \mathcal{F}$. For the inductive step: assuming $B_n \in \mathcal{F}$,
we write $B_{n+1} = B_n \cup A_{n+1}$. Now:
\begin{equation}
    B_n \cup A_{n+1} = (B_n^c \cap A_{n+1}^c)^c.
\end{equation}
Since $\mathcal{F}$ is closed under complement (shown above) and intersection
($\pi$-system), $B_n^c \cap A_{n+1}^c \in \mathcal{F}$, hence
$(B_n^c \cap A_{n+1}^c)^c \in \mathcal{F}$.

So $B_n$ is an increasing sequence in $\mathcal{F}$, and $\bigcup_{n=1}^\infty A_n = \bigcup_{n=1}^\infty B_n \in \mathcal{F}$ by the $\lambda$-system property.
\end{proof}

\begin{keyresult}[Dynkin's $\pi$--$\lambda$ Theorem]
\begin{theorem}[$\pi$--$\lambda$ Theorem]
\label{thm:pi_lambda}
If $\mathcal{P}$ is a $\pi$-system and $\mathcal{L}$ is a $\lambda$-system with
$\mathcal{P} \subseteq \mathcal{L}$, then $\sigma(\mathcal{P}) \subseteq \mathcal{L}$.
\end{theorem}
\end{keyresult}

\begin{proof}
Let $\ell(\mathcal{P})$ denote the smallest $\lambda$-system containing $\mathcal{P}$
(this exists by the same intersection argument as for generated $\sigma$-algebras:
$\powerset(\Omega)$ is a $\lambda$-system, and the intersection of $\lambda$-systems
is a $\lambda$-system). Since $\mathcal{L}$ is a $\lambda$-system containing $\mathcal{P}$,
we have $\ell(\mathcal{P}) \subseteq \mathcal{L}$.

It suffices to show $\ell(\mathcal{P})$ is a $\sigma$-algebra (for then $\sigma(\mathcal{P}) \subseteq \ell(\mathcal{P}) \subseteq \mathcal{L}$). By Lemma~\ref{lem:pi_lambda_sigma}, it suffices to show $\ell(\mathcal{P})$ is a $\pi$-system.

\textbf{Step 1:} For each $B \in \ell(\mathcal{P})$, define
$\mathcal{G}_B = \{A \in \ell(\mathcal{P}) : A \cap B \in \ell(\mathcal{P})\}$.
We claim $\mathcal{G}_B$ is a $\lambda$-system:
\begin{itemize}[nosep]
    \item $\Omega \in \mathcal{G}_B$ since $\Omega \cap B = B \in \ell(\mathcal{P})$.
    \item If $A_1 \subseteq A_2$ with $A_1, A_2 \in \mathcal{G}_B$, then $(A_2 \setminus A_1) \cap B = (A_2 \cap B) \setminus (A_1 \cap B)$, and $A_1 \cap B \subseteq A_2 \cap B$ with both in $\ell(\mathcal{P})$, so $(A_2 \setminus A_1) \cap B \in \ell(\mathcal{P})$.
    \item If $A_n \uparrow$, then $A_n \cap B \uparrow$, so $(\bigcup A_n) \cap B = \bigcup(A_n \cap B) \in \ell(\mathcal{P})$.
\end{itemize}

\textbf{Step 2:} For $B \in \mathcal{P}$: since $\mathcal{P}$ is a $\pi$-system,
$A \cap B \in \mathcal{P} \subseteq \ell(\mathcal{P})$ for all $A \in \mathcal{P}$.
So $\mathcal{P} \subseteq \mathcal{G}_B$. Since $\mathcal{G}_B$ is a $\lambda$-system,
$\ell(\mathcal{P}) \subseteq \mathcal{G}_B$. This means: for all $B \in \mathcal{P}$
and $A \in \ell(\mathcal{P})$, $A \cap B \in \ell(\mathcal{P})$.

\textbf{Step 3:} Now fix $A \in \ell(\mathcal{P})$ and consider
$\mathcal{G}_A = \{B \in \ell(\mathcal{P}) : A \cap B \in \ell(\mathcal{P})\}$.
By Step 2, $\mathcal{P} \subseteq \mathcal{G}_A$ (since for $B \in \mathcal{P}$,
$A \cap B \in \ell(\mathcal{P})$ by Step 2). Since $\mathcal{G}_A$ is a
$\lambda$-system (by Step 1 with $A$ and $B$ swapped), $\ell(\mathcal{P}) \subseteq \mathcal{G}_A$.

This means: for all $A, B \in \ell(\mathcal{P})$, $A \cap B \in \ell(\mathcal{P})$.
So $\ell(\mathcal{P})$ is a $\pi$-system. By Lemma~\ref{lem:pi_lambda_sigma},
$\ell(\mathcal{P})$ is a $\sigma$-algebra, so $\sigma(\mathcal{P}) \subseteq \ell(\mathcal{P}) \subseteq \mathcal{L}$.
\end{proof}

\begin{corollary}[Uniqueness of Measures]
\label{cor:uniqueness}
If two probability measures $\Prob_1$ and $\Prob_2$ on $(\Omega, \sigma(\mathcal{P}))$
agree on a $\pi$-system $\mathcal{P}$, then $\Prob_1 = \Prob_2$ on
$\sigma(\mathcal{P})$.
\end{corollary}

\begin{proof}
The collection $\mathcal{L} = \{A \in \sigma(\mathcal{P}) : \Prob_1(A) = \Prob_2(A)\}$
is a $\lambda$-system containing $\mathcal{P}$. By the $\pi$--$\lambda$ theorem,
$\sigma(\mathcal{P}) \subseteq \mathcal{L}$.

We verify $\mathcal{L}$ is a $\lambda$-system:
\begin{itemize}[nosep]
    \item $\Omega \in \mathcal{L}$ since $\Prob_1(\Omega) = \Prob_2(\Omega) = 1$.
    \item If $A \subseteq B$ with $A, B \in \mathcal{L}$, then $\Prob_i(B \setminus A) = \Prob_i(B) - \Prob_i(A)$ for $i = 1, 2$, and these are equal.
    \item If $A_n \uparrow A$ with $A_n \in \mathcal{L}$, then $\Prob_i(A) = \lim \Prob_i(A_n)$ by continuity from below, and these limits are equal. \qedhere
\end{itemize}
\end{proof}

\begin{rlconnection}[Connection to RL]
Corollary~\ref{cor:uniqueness} is why the distribution of a random variable
is uniquely determined by its CDF $F_X(x) = \Prob(X \leq x)$: the half-lines
$(-\infty, x]$ form a $\pi$-system generating $\Borel(\R)$, so any two measures
agreeing on all $(-\infty, x]$ must agree on all Borel sets. In RL, this means
the distribution of the return $G_t$ is fully characterized by its CDF --- there
is no ``hidden'' probabilistic information.
\end{rlconnection}

%% ============================================================
%% SECTION 3: MEASURES
%% ============================================================
\section{Measures}
\label{sec:measures}

\subsection{Definition and Fundamental Properties}
\label{subsec:measure_def}

\begin{definition}[Measure]
\label{def:measure}
A \emph{measure} on a measurable space $(\Omega, \mathcal{F})$ is a function
$\mu: \mathcal{F} \to [0, \infty]$ satisfying:
\begin{enumerate}[nosep,label=(\roman*)]
    \item $\mu(\emptyset) = 0$.
    \item \textbf{Countable additivity:} If $A_1, A_2, \ldots \in \mathcal{F}$ are
    pairwise disjoint, then:
    \begin{equation}
        \mu\!\left(\bigsqcup_{n=1}^\infty A_n\right) = \sum_{n=1}^\infty \mu(A_n).
        \label{eq:countable_additivity}
    \end{equation}
\end{enumerate}
The triple $(\Omega, \mathcal{F}, \mu)$ is a \emph{measure space}. If
$\mu(\Omega) = 1$, then $\mu$ is a \emph{probability measure} and we write
$\Prob$ instead of $\mu$.
\end{definition}

\begin{example}[Standard Measures]
\label{ex:measures}
\begin{enumerate}[nosep,label=(\alph*)]
    \item \textbf{Counting measure:} $\mu(A) = |A|$ (number of elements) on any countable $\Omega$.
    \item \textbf{Dirac measure:} $\delta_x(A) = \indicator_A(x) = \begin{cases} 1 & \text{if } x \in A \\ 0 & \text{if } x \notin A \end{cases}$.
    \item \textbf{Lebesgue measure:} $\lambda((a,b]) = b - a$ on $(\R, \Borel(\R))$. (Existence proved in Section~\ref{sec:construction}.)
    \item \textbf{Probability on a finite set:} $\Prob(\{s\}) = p(s)$ with $\sum_{s \in \mathcal{S}} p(s) = 1$.
\end{enumerate}
\end{example}

\begin{theorem}[Properties of Measures]
\label{thm:measure_props}
Let $(\Omega, \mathcal{F}, \mu)$ be a measure space. For $A, B, A_n \in \mathcal{F}$:
\begin{enumerate}[nosep,label=(\alph*)]
    \item \textbf{Monotonicity:} $A \subseteq B \implies \mu(A) \leq \mu(B)$.
    \item \textbf{Subtractivity:} $A \subseteq B$ and $\mu(A) < \infty \implies \mu(B \setminus A) = \mu(B) - \mu(A)$.
    \item \textbf{Sub-additivity:} $\mu\!\left(\bigcup_{n=1}^\infty A_n\right) \leq \sum_{n=1}^\infty \mu(A_n)$.
    \item \textbf{Continuity from below:} $A_n \uparrow A \implies \mu(A) = \lim_n \mu(A_n)$.
    \item \textbf{Continuity from above:} $A_n \downarrow A$ and $\mu(A_1) < \infty \implies \mu(A) = \lim_n \mu(A_n)$.
\end{enumerate}
\end{theorem}

\begin{proof}
\textbf{(a)} $B = A \sqcup (B \setminus A)$, so $\mu(B) = \mu(A) + \mu(B \setminus A) \geq \mu(A)$.

\textbf{(b)} From $\mu(B) = \mu(A) + \mu(B \setminus A)$ and $\mu(A) < \infty$, subtract $\mu(A)$.

\textbf{(c)} Define $B_1 = A_1$, $B_n = A_n \setminus \bigcup_{k=1}^{n-1} A_k$ for $n \geq 2$. The $B_n$ are disjoint, $\bigcup B_n = \bigcup A_n$, and $B_n \subseteq A_n$. So:
$\mu(\bigcup A_n) = \mu(\bigsqcup B_n) = \sum \mu(B_n) \leq \sum \mu(A_n)$.

\textbf{(d)} Let $A_0 = \emptyset$. The sets $C_n = A_n \setminus A_{n-1}$ are disjoint and $A_n = \bigsqcup_{k=1}^n C_k$. Also $A = \bigsqcup_{k=1}^\infty C_k$. By countable additivity:
\begin{equation}
    \mu(A) = \sum_{k=1}^\infty \mu(C_k) = \lim_{n \to \infty} \sum_{k=1}^n \mu(C_k) = \lim_{n \to \infty} \mu(A_n).
\end{equation}

\textbf{(e)} Define $B_n = A_1 \setminus A_n$. Then $B_n \uparrow A_1 \setminus A$. By (d):
$\mu(A_1 \setminus A) = \lim \mu(B_n) = \lim (\mu(A_1) - \mu(A_n))$, using (b) (valid since $\mu(A_n) \leq \mu(A_1) < \infty$). So $\mu(A_1) - \mu(A) = \mu(A_1) - \lim \mu(A_n)$, giving $\mu(A) = \lim \mu(A_n)$.
\end{proof}

\begin{warning}[Finiteness Condition in (e)]
The condition $\mu(A_1) < \infty$ in continuity from above is essential. Consider
counting measure on $\N$ with $A_n = \{n, n+1, n+2, \ldots\}$. Then $A_n \downarrow \emptyset$
but $\mu(A_n) = \infty$ for all $n$, while $\mu(\emptyset) = 0$. For probability
measures, this condition is always satisfied since $\Prob(A_1) \leq 1 < \infty$.
\end{warning}

\subsection{Null Sets and Complete Measures}
\label{subsec:null_sets}

\begin{definition}[Null Set, Complete Measure]
\label{def:null_complete}
A set $N \in \mathcal{F}$ is a \emph{$\mu$-null set} (or \emph{null set}) if $\mu(N) = 0$.
A measure $\mu$ is \emph{complete} if whenever $N$ is $\mu$-null and $A \subseteq N$,
then $A \in \mathcal{F}$ (and hence $\mu(A) = 0$).
\end{definition}

\begin{proposition}[Completion of a Measure]
\label{prop:completion}
Every measure space $(\Omega, \mathcal{F}, \mu)$ can be completed: there exists
a smallest $\sigma$-algebra $\overline{\mathcal{F}} \supseteq \mathcal{F}$ and
a measure $\bar{\mu}$ extending $\mu$ such that $(\Omega, \overline{\mathcal{F}}, \bar{\mu})$
is complete. Explicitly:
\begin{equation}
    \overline{\mathcal{F}} = \{A \cup N' : A \in \mathcal{F},\; N' \subseteq N \text{ for some $\mu$-null } N \in \mathcal{F}\},
\end{equation}
and $\bar{\mu}(A \cup N') = \mu(A)$.
\end{proposition}

\begin{proof}
We verify that $\overline{\mathcal{F}}$ is a $\sigma$-algebra. The key step is
closure under complement: if $E = A \cup N'$ with $N' \subseteq N$ and $\mu(N) = 0$,
then $E^c = A^c \cap (N')^c = (A^c \cap N^c) \cup (A^c \cap N \cap (N')^c)$.
The first part is in $\mathcal{F}$, and the second is a subset of $N$ (a null set),
so $E^c \in \overline{\mathcal{F}}$.

For well-definedness of $\bar{\mu}$: if $A_1 \cup N_1' = A_2 \cup N_2'$, then
$A_1 \triangle A_2 \subseteq N_1 \cup N_2$ (a null set), so
$\mu(A_1 \triangle A_2) = 0$, hence $\mu(A_1) = \mu(A_2)$.
\end{proof}

%% ============================================================
%% SECTION 4: CONSTRUCTION OF MEASURES
%% ============================================================
\section{Construction of Measures: Carath\'{e}odory's Theorem}
\label{sec:construction}

How do we actually \emph{build} measures? We cannot simply assign
$\mu(A)$ for every $A \in \Borel(\R)$ individually --- the Borel $\sigma$-algebra
is uncountably large. Instead, we specify $\mu$ on a small, manageable collection
and extend it.

\subsection{Algebras and Pre-Measures}
\label{subsec:algebras}

\begin{definition}[Algebra of Sets]
\label{def:algebra}
A collection $\mathcal{A} \subseteq \powerset(\Omega)$ is an \emph{algebra}
(or \emph{field}) if:
\begin{enumerate}[nosep,label=(\roman*)]
    \item $\Omega \in \mathcal{A}$.
    \item $A \in \mathcal{A} \implies A^c \in \mathcal{A}$.
    \item $A, B \in \mathcal{A} \implies A \cup B \in \mathcal{A}$.
    \hfill(closed under \emph{finite} union)
\end{enumerate}
\end{definition}

\begin{remark}
An algebra is like a $\sigma$-algebra but with only finite (not countable) closure.
Every $\sigma$-algebra is an algebra, but not conversely.
\end{remark}

\begin{example}
\label{ex:algebra}
On $\R$, the collection $\mathcal{A}$ of all finite unions of half-open intervals
$(a, b]$ (including $\emptyset$ and $\R = (-\infty, \infty)$) is an algebra
but not a $\sigma$-algebra. For instance, $\bigcup_{n=1}^\infty (0, 1 - 1/n] = (0, 1)$
is not a finite union of half-open intervals, so $(0,1) \notin \mathcal{A}$.
\end{example}

\begin{definition}[Pre-Measure]
\label{def:premeasure}
A \emph{pre-measure} on an algebra $\mathcal{A}$ is a function
$\mu_0: \mathcal{A} \to [0, \infty]$ satisfying:
\begin{enumerate}[nosep,label=(\roman*)]
    \item $\mu_0(\emptyset) = 0$.
    \item If $A_1, A_2, \ldots \in \mathcal{A}$ are disjoint and
    $\bigsqcup_{n=1}^\infty A_n \in \mathcal{A}$, then
    $\mu_0(\bigsqcup A_n) = \sum \mu_0(A_n)$.
\end{enumerate}
\end{definition}

\subsection{Outer Measures}
\label{subsec:outer}

\begin{definition}[Outer Measure]
\label{def:outer}
An \emph{outer measure} on $\Omega$ is a function
$\mu^*: \powerset(\Omega) \to [0, \infty]$ satisfying:
\begin{enumerate}[nosep,label=(\roman*)]
    \item $\mu^*(\emptyset) = 0$.
    \item \textbf{Monotonicity:} $A \subseteq B \implies \mu^*(A) \leq \mu^*(B)$.
    \item \textbf{Countable sub-additivity:} $\mu^*(\bigcup_{n=1}^\infty A_n) \leq \sum_{n=1}^\infty \mu^*(A_n)$.
\end{enumerate}
\end{definition}

\begin{proposition}[Pre-Measure Induces Outer Measure]
\label{prop:induced_outer}
Given a pre-measure $\mu_0$ on an algebra $\mathcal{A}$, define for any
$E \subseteq \Omega$:
\begin{equation}
    \mu^*(E) = \inf\left\{\sum_{n=1}^\infty \mu_0(A_n) : A_n \in \mathcal{A},\; E \subseteq \bigcup_{n=1}^\infty A_n\right\}.
    \label{eq:induced_outer}
\end{equation}
Then $\mu^*$ is an outer measure, and $\mu^*(A) = \mu_0(A)$ for all $A \in \mathcal{A}$.
\end{proposition}

\begin{proof}
\textit{(i)} $\mu^*(\emptyset) = 0$: cover $\emptyset$ by $A_n = \emptyset$ for all $n$.

\textit{(ii)} If $A \subseteq B$, every cover of $B$ is also a cover of $A$, so the infimum over covers of $A$ is at most the infimum over covers of $B$.

\textit{(iii)} Let $E = \bigcup_n E_n$. For each $\varepsilon > 0$ and each $n$, choose
a cover $\{A_{n,k}\}_{k} \subseteq \mathcal{A}$ with $E_n \subseteq \bigcup_k A_{n,k}$
and $\sum_k \mu_0(A_{n,k}) \leq \mu^*(E_n) + \varepsilon/2^n$. Then
$\{A_{n,k}\}_{n,k}$ covers $E$, so:
\begin{equation}
    \mu^*(E) \leq \sum_n \sum_k \mu_0(A_{n,k}) \leq \sum_n (\mu^*(E_n) + \varepsilon/2^n) = \sum_n \mu^*(E_n) + \varepsilon.
\end{equation}
Since $\varepsilon$ is arbitrary, $\mu^*(E) \leq \sum_n \mu^*(E_n)$.

\textit{$\mu^*$ extends $\mu_0$:} For $A \in \mathcal{A}$, the single-set cover gives
$\mu^*(A) \leq \mu_0(A)$. For the reverse: if $A \subseteq \bigcup_n A_n$ with
$A_n \in \mathcal{A}$, define $B_n = A \cap A_n \cap (\bigcup_{k<n} A_k)^c \in \mathcal{A}$
(since $\mathcal{A}$ is an algebra). Then $A = \bigsqcup B_n$ and $B_n \subseteq A_n$.
By the pre-measure properties: $\mu_0(A) = \sum \mu_0(B_n) \leq \sum \mu_0(A_n)$.
Taking the infimum: $\mu_0(A) \leq \mu^*(A)$.
\end{proof}

\subsection{Carath\'{e}odory's Extension Theorem}
\label{subsec:caratheodory}

\begin{definition}[Carath\'{e}odory Measurability]
\label{def:cara_measurable}
Given an outer measure $\mu^*$ on $\Omega$, a set $A \subseteq \Omega$ is
\emph{$\mu^*$-measurable} (or \emph{Carath\'{e}odory measurable}) if for every
$E \subseteq \Omega$:
\begin{equation}
    \mu^*(E) = \mu^*(E \cap A) + \mu^*(E \cap A^c).
    \label{eq:cara_condition}
\end{equation}
\end{definition}

\begin{remark}
By sub-additivity, $\mu^*(E) \leq \mu^*(E \cap A) + \mu^*(E \cap A^c)$ always holds.
So the content of~\eqref{eq:cara_condition} is the reverse inequality: $A$ ``splits''
every test set $E$ additively. This is a strong condition --- it must hold for
\emph{all} $E \subseteq \Omega$, not just sets in some $\sigma$-algebra.
\end{remark}

\begin{keyresult}[Carath\'{e}odory's Extension Theorem]
\begin{theorem}[Carath\'{e}odory]
\label{thm:caratheodory}
Let $\mu^*$ be an outer measure on $\Omega$. Then:
\begin{enumerate}[nosep,label=(\alph*)]
    \item The collection $\mathcal{M} = \{A \subseteq \Omega : A \text{ is $\mu^*$-measurable}\}$ is a $\sigma$-algebra.
    \item The restriction $\mu = \mu^*|_{\mathcal{M}}$ is a complete measure on $(\Omega, \mathcal{M})$.
\end{enumerate}
Consequently, if $\mu_0$ is a pre-measure on an algebra $\mathcal{A}$ and $\mu^*$ is
the induced outer measure~\eqref{eq:induced_outer}, then $\mathcal{A} \subseteq \mathcal{M}$
and $\mu|_{\mathcal{A}} = \mu_0$. If additionally $\mu_0$ is $\sigma$-finite, then
$\mu$ is the unique extension of $\mu_0$ to $\sigma(\mathcal{A})$.
\end{theorem}
\end{keyresult}

\begin{proof}
\textbf{Part (a): $\mathcal{M}$ is a $\sigma$-algebra.}

\textit{Step 1: $\Omega \in \mathcal{M}$.} For any $E$: $\mu^*(E \cap \Omega) + \mu^*(E \cap \emptyset) = \mu^*(E) + 0 = \mu^*(E)$. \checkmark

\textit{Step 2: Closure under complement.} The condition~\eqref{eq:cara_condition} is symmetric in $A$ and $A^c$, so $A \in \mathcal{M} \Leftrightarrow A^c \in \mathcal{M}$. \checkmark

\textit{Step 3: Closure under finite union.} Let $A, B \in \mathcal{M}$. For any test set $E$:
\begin{align}
    \mu^*(E) &= \mu^*(E \cap A) + \mu^*(E \cap A^c) \tag{$A$ measurable}\\
    &= \mu^*(E \cap A) + \mu^*(E \cap A^c \cap B) + \mu^*(E \cap A^c \cap B^c) \tag{$B$ measurable, applied to $E \cap A^c$}
\end{align}
We need $\mu^*(E) = \mu^*(E \cap (A \cup B)) + \mu^*(E \cap (A \cup B)^c)$.

Note $(A \cup B)^c = A^c \cap B^c$, so the third term is correct. For the first:
$E \cap (A \cup B) = (E \cap A) \cup (E \cap A^c \cap B)$, so by sub-additivity:
$\mu^*(E \cap (A \cup B)) \leq \mu^*(E \cap A) + \mu^*(E \cap A^c \cap B)$.

Combining: $\mu^*(E) \geq \mu^*(E \cap (A \cup B)) + \mu^*(E \cap (A \cup B)^c)$.
The reverse holds by sub-additivity. So $A \cup B \in \mathcal{M}$. \checkmark

\textit{Step 4: Closure under countable disjoint union.} Let $A_1, A_2, \ldots \in \mathcal{M}$ be disjoint. Set $S_n = \bigsqcup_{k=1}^n A_k$ and $S = \bigsqcup_{k=1}^\infty A_k$.

By induction (using Step 3 and closure under complement, hence closure under intersection), $S_n \in \mathcal{M}$, and for any $E$:
\begin{equation}
    \mu^*(E \cap S_n) = \sum_{k=1}^n \mu^*(E \cap A_k).
    \label{eq:cara_finite}
\end{equation}
This is proved by induction: $\mu^*(E \cap S_n) = \mu^*(E \cap S_n \cap A_n) + \mu^*(E \cap S_n \cap A_n^c) = \mu^*(E \cap A_n) + \mu^*(E \cap S_{n-1})$.

Since $S_n \in \mathcal{M}$: $\mu^*(E) = \mu^*(E \cap S_n) + \mu^*(E \cap S_n^c)$.
Since $S_n \subseteq S$, $S^c \subseteq S_n^c$, so $\mu^*(E \cap S_n^c) \geq \mu^*(E \cap S^c)$. Using~\eqref{eq:cara_finite}:
\begin{equation}
    \mu^*(E) \geq \sum_{k=1}^n \mu^*(E \cap A_k) + \mu^*(E \cap S^c).
\end{equation}
Letting $n \to \infty$:
\begin{equation}
    \mu^*(E) \geq \sum_{k=1}^\infty \mu^*(E \cap A_k) + \mu^*(E \cap S^c) \geq \mu^*(E \cap S) + \mu^*(E \cap S^c),
\end{equation}
where the last step uses sub-additivity $\mu^*(E \cap S) \leq \sum \mu^*(E \cap A_k)$.
The reverse inequality holds by sub-additivity. So $S \in \mathcal{M}$. \checkmark

\textit{Step 5: General countable unions.} For arbitrary (not necessarily disjoint)
$A_n \in \mathcal{M}$, write $\bigcup A_n$ as a disjoint union using $B_n = A_n \setminus \bigcup_{k<n} A_k \in \mathcal{M}$, and apply Step 4.

\textbf{Part (b): $\mu$ is a complete measure.}

Countable additivity follows from~\eqref{eq:cara_finite} with $E = S$:
$\mu(S) = \mu^*(S) = \sum \mu^*(A_k) = \sum \mu(A_k)$.

Completeness: if $\mu^*(N) = 0$ and $A \subseteq N$, then for any $E$:
$\mu^*(E \cap A) \leq \mu^*(A) \leq \mu^*(N) = 0$, so
$\mu^*(E \cap A) + \mu^*(E \cap A^c) = \mu^*(E \cap A^c) \leq \mu^*(E)$.
The reverse holds by monotonicity ($E \cap A^c \subseteq E$). Wait, actually we need
$\mu^*(E \cap A^c) \geq \mu^*(E) - \mu^*(E \cap A) = \mu^*(E)$ and $\mu^*(E \cap A^c) \leq \mu^*(E)$. These give equality. So $A \in \mathcal{M}$.

\textbf{Uniqueness under $\sigma$-finiteness:} follows from Corollary~\ref{cor:uniqueness}
applied to the $\pi$-system $\mathcal{A}$, after reducing to the finite-measure case
using the $\sigma$-finite decomposition.
\end{proof}

\subsection{Application: Lebesgue Measure}
\label{subsec:lebesgue}

\begin{theorem}[Existence of Lebesgue Measure]
\label{thm:lebesgue}
There exists a unique measure $\lambda$ on $(\R, \Borel(\R))$ such that
$\lambda((a, b]) = b - a$ for all $a < b$. This measure is $\sigma$-finite and
translation-invariant: $\lambda(A + x) = \lambda(A)$ for all $A \in \Borel(\R)$
and $x \in \R$.
\end{theorem}

\begin{proof}[Proof sketch]
Let $\mathcal{A}$ be the algebra of finite disjoint unions of half-open intervals
$(a_i, b_i]$. Define $\mu_0(\bigsqcup_i (a_i, b_i]) = \sum_i (b_i - a_i)$. One
verifies $\mu_0$ is a pre-measure on $\mathcal{A}$ (the key step is showing
countable additivity, which uses the Heine--Borel theorem: if $(a, b]$ is covered
by countably many $(a_n, b_n]$, then $b - a \leq \sum (b_n - a_n)$).

By Carath\'{e}odory's theorem (Theorem~\ref{thm:caratheodory}), $\mu_0$ extends
uniquely to $\sigma(\mathcal{A}) = \Borel(\R)$ since $\mu_0$ is $\sigma$-finite
($\R = \bigcup_{n=1}^\infty (-n, n]$ with $\mu_0((-n, n]) = 2n < \infty$).
\end{proof}

\begin{theorem}[Non-Measurable Sets (Vitali)]
\label{thm:vitali}
There exists a set $V \subseteq [0, 1]$ that is not Lebesgue measurable.
\end{theorem}

\begin{proof}
Define an equivalence relation on $[0,1)$: $x \sim y$ iff $x - y \in \Q$.
By the Axiom of Choice, select one representative from each equivalence class;
call this set $V$.

For $q \in \Q \cap [0, 1)$, define $V_q = \{x + q \mod 1 : x \in V\}$
(shift $V$ by $q$ modulo 1). The $V_q$ are disjoint (if $v_1 + q_1 \equiv v_2 + q_2$
then $v_1 - v_2 \in \Q$, so $v_1 \sim v_2$, so $v_1 = v_2$ and $q_1 = q_2$).
Also $[0, 1) = \bigsqcup_{q \in \Q \cap [0,1)} V_q$ (every $x$ is equivalent to
some $v \in V$, and $x = v + q$ for some rational $q$).

If $V$ were measurable, then each $V_q$ would have the same measure (by
translation invariance of Lebesgue measure modulo 1). Call this $c$.

By countable additivity: $1 = \lambda([0,1)) = \sum_{q} c$. But $\sum c$ is either
$0$ (if $c = 0$) or $\infty$ (if $c > 0$) --- never $1$. Contradiction.
\end{proof}

\begin{remark}
Vitali's theorem shows why we cannot define Lebesgue measure on $\powerset(\R)$:
not all subsets of $\R$ are measurable. This is the fundamental reason for
restricting to $\sigma$-algebras. In probability, this subtlety is rarely
encountered (finite MDPs use $\powerset(\mathcal{S})$), but it justifies
the formalism and is important for continuous state/action spaces.
\end{remark}

%% ============================================================
%% SECTION 5: MEASURABLE FUNCTIONS
%% ============================================================
\section{Measurable Functions}
\label{sec:measurable}

\subsection{Definition and Characterization}
\label{subsec:meas_def}

\begin{definition}[Measurable Function]
\label{def:measurable_fn}
Let $(\Omega, \mathcal{F})$ and $(S, \mathcal{S})$ be measurable spaces. A function
$f: \Omega \to S$ is \emph{$(\mathcal{F}, \mathcal{S})$-measurable} (or simply
\emph{measurable}) if:
\begin{equation}
    f^{-1}(B) := \{\omega \in \Omega : f(\omega) \in B\} \in \mathcal{F} \quad \text{for all } B \in \mathcal{S}.
    \label{eq:measurable_fn}
\end{equation}
When $S = \R$ and $\mathcal{S} = \Borel(\R)$, we say $f$ is \emph{Borel measurable}
(or just \emph{measurable}).
\end{definition}

\begin{theorem}[Generator Criterion for Measurability]
\label{thm:generator_criterion}
Let $\mathcal{S} = \sigma(\mathcal{C})$ for some collection $\mathcal{C}$. Then
$f: (\Omega, \mathcal{F}) \to (S, \mathcal{S})$ is measurable if and only if:
\begin{equation}
    f^{-1}(C) \in \mathcal{F} \quad \text{for all } C \in \mathcal{C}.
\end{equation}
\end{theorem}

\begin{proof}
($\Rightarrow$) Trivial: $\mathcal{C} \subseteq \mathcal{S}$.

($\Leftarrow$) Define $\mathcal{G} = \{B \in \mathcal{S} : f^{-1}(B) \in \mathcal{F}\}$. We show $\mathcal{G}$ is a $\sigma$-algebra:
\begin{itemize}[nosep]
    \item $S \in \mathcal{G}$: $f^{-1}(S) = \Omega \in \mathcal{F}$.
    \item $B \in \mathcal{G} \implies B^c \in \mathcal{G}$: $f^{-1}(B^c) = (f^{-1}(B))^c \in \mathcal{F}$.
    \item $B_n \in \mathcal{G} \implies \bigcup B_n \in \mathcal{G}$: $f^{-1}(\bigcup B_n) = \bigcup f^{-1}(B_n) \in \mathcal{F}$.
\end{itemize}
By hypothesis, $\mathcal{C} \subseteq \mathcal{G}$. Since $\mathcal{G}$ is a $\sigma$-algebra, $\sigma(\mathcal{C}) = \mathcal{S} \subseteq \mathcal{G}$, so $\mathcal{G} = \mathcal{S}$.
\end{proof}

\begin{corollary}[Standard Criterion for Real-Valued Measurability]
\label{cor:real_measurable}
$f: (\Omega, \mathcal{F}) \to (\R, \Borel(\R))$ is measurable if and only if
$\{f \leq x\} := \{\omega : f(\omega) \leq x\} \in \mathcal{F}$ for all $x \in \R$.
\end{corollary}

\begin{proof}
Apply Theorem~\ref{thm:generator_criterion} with $\mathcal{C} = \{(-\infty, x] : x \in \R\}$, which generates $\Borel(\R)$ by Proposition~\ref{prop:borel_generators}(d).
\end{proof}

\begin{rlconnection}[Connection to RL]
A \emph{random variable} is precisely a measurable function $X: (\Omega, \mathcal{F}) \to (\R, \Borel(\R))$. In an MDP, the state $S_t$, action $A_t$, and reward $R_{t+1}$ are all random variables. The value function $V^\pi(s) = \E[G_t \mid S_t = s]$ is a measurable function of $s$. Corollary~\ref{cor:real_measurable} says that being a random variable means precisely that events like $\{X \leq x\}$ are ``observable'' (in $\mathcal{F}$) --- which is why we can compute their probabilities.
\end{rlconnection}

\subsection{Algebra of Measurable Functions}
\label{subsec:meas_algebra}

\begin{theorem}[Closure Properties]
\label{thm:meas_closure}
Let $f, g: (\Omega, \mathcal{F}) \to (\R, \Borel(\R))$ be measurable and $c \in \R$.
Then the following are measurable:
\begin{enumerate}[nosep,label=(\alph*)]
    \item $cf$, \quad $f + g$, \quad $fg$, \quad $|f|$, \quad $\max(f, g)$, \quad $\min(f, g)$.
    \item $f/g$ on $\{g \neq 0\}$.
    \item $f^+ = \max(f, 0)$ and $f^- = \max(-f, 0)$.
\end{enumerate}
\end{theorem}

\begin{proof}
We prove $f + g$ is measurable; the others are similar.

\textbf{Key idea:} $\{f + g < c\} = \bigcup_{q \in \Q} (\{f < q\} \cap \{g < c - q\})$.

\textit{Why?} If $f(\omega) + g(\omega) < c$, then $f(\omega) < c - g(\omega)$. By density of $\Q$, there exists $q \in \Q$ with $f(\omega) < q < c - g(\omega)$, so $\omega \in \{f < q\} \cap \{g < c - q\}$.

Conversely, if $f(\omega) < q$ and $g(\omega) < c - q$, then $f(\omega) + g(\omega) < c$.

Since each $\{f < q\}$ and $\{g < c - q\}$ is in $\mathcal{F}$ (measurability of $f$ and $g$), and the union is countable (over $q \in \Q$), we conclude $\{f + g < c\} \in \mathcal{F}$.

Since $\{f + g < c\} \in \mathcal{F}$ for all $c$ and $\{f + g \leq c\} = \bigcap_{n} \{f + g < c + 1/n\}$, we get $\{f + g \leq c\} \in \mathcal{F}$.
\end{proof}

\begin{theorem}[Limits of Measurable Functions]
\label{thm:meas_limits}
If $f_1, f_2, \ldots$ are measurable, then the following are measurable:
\begin{equation}
    \sup_n f_n, \quad \inf_n f_n, \quad \limsup_n f_n, \quad \liminf_n f_n.
\end{equation}
In particular, if $f_n \to f$ pointwise (or a.s.), then $f$ is measurable.
\end{theorem}

\begin{proof}
$\{\sup_n f_n \leq c\} = \bigcap_{n=1}^\infty \{f_n \leq c\} \in \mathcal{F}$ (countable intersection of measurable sets).

$\inf_n f_n = -\sup_n(-f_n)$ is measurable.

$\limsup_n f_n = \inf_n \sup_{k \geq n} f_k$ and $\liminf_n f_n = \sup_n \inf_{k \geq n} f_k$ are measurable by composing the above.

If $f_n \to f$ pointwise, then $f = \lim_n f_n = \limsup_n f_n = \liminf_n f_n$.
\end{proof}

\subsection{Simple Functions and Approximation}
\label{subsec:simple}

\begin{definition}[Simple Function]
\label{def:simple}
A \emph{simple function} is a measurable function $\varphi: \Omega \to \R$ taking
finitely many values. Every simple function can be written as:
\begin{equation}
    \varphi = \sum_{i=1}^n a_i \indicator_{A_i},
    \label{eq:simple}
\end{equation}
where $a_1, \ldots, a_n \in \R$ are distinct, $A_i = \varphi^{-1}(\{a_i\}) \in \mathcal{F}$,
and $\Omega = \bigsqcup_{i=1}^n A_i$. This is the \emph{standard representation}.
\end{definition}

\begin{keyresult}[Simple Function Approximation]
\begin{theorem}[Approximation Theorem]
\label{thm:approximation}
Let $f: (\Omega, \mathcal{F}) \to [0, \infty]$ be measurable. Then there exist
simple functions $0 \leq \varphi_1 \leq \varphi_2 \leq \cdots$ with $\varphi_n \uparrow f$
pointwise. If $f$ is bounded, the convergence is uniform.
\end{theorem}
\end{keyresult}

\begin{proof}
For each $n \geq 1$, define:
\begin{equation}
    \varphi_n(\omega) = \begin{cases}
        \frac{k-1}{2^n} & \text{if } \frac{k-1}{2^n} \leq f(\omega) < \frac{k}{2^n}, \quad k = 1, 2, \ldots, n \cdot 2^n,\\[4pt]
        n & \text{if } f(\omega) \geq n.
    \end{cases}
    \label{eq:approx_construction}
\end{equation}
In other words, $\varphi_n$ rounds $f$ down to the nearest multiple of $1/2^n$,
capped at $n$.

\textit{Measurability:} $\varphi_n$ takes finitely many values, and each level set
$\varphi_n^{-1}(\{k/2^n\}) = \{k/2^n \leq f < (k+1)/2^n\} \in \mathcal{F}$.

\textit{Monotonicity:} $\varphi_{n+1}$ uses a finer grid ($1/2^{n+1}$ vs $1/2^n$)
and a higher cap ($n+1$ vs $n$), so $\varphi_{n+1} \geq \varphi_n$.

\textit{Convergence:} If $f(\omega) < \infty$, then for $n > f(\omega)$:
$0 \leq f(\omega) - \varphi_n(\omega) < 1/2^n \to 0$.
If $f(\omega) = \infty$, then $\varphi_n(\omega) = n \to \infty = f(\omega)$.

\textit{Uniform convergence when bounded:} If $f \leq M$, then for $n > M$,
$|f - \varphi_n| < 1/2^n$ everywhere.
\end{proof}

\begin{rlconnection}[Connection to RL]
The approximation theorem is the foundation of Lebesgue integration (next lesson).
The Lebesgue integral of a non-negative function is defined as
$\int f\,d\mu = \sup\{\int \varphi\,d\mu : \varphi \text{ simple}, \varphi \leq f\}$,
where simple function integrals are just weighted sums. In RL, this underlies
the expectation $\E[G_t \mid S_t = s]$: the return $G_t$ is approximated by
simple functions, and the expectation is the limit of their integrals.
\end{rlconnection}

\subsection{$\sigma$-Algebra Generated by a Random Variable}
\label{subsec:sigma_rv}

\begin{definition}[$\sigma(X)$]
\label{def:sigma_x}
For a measurable function $X: (\Omega, \mathcal{F}) \to (S, \mathcal{S})$, the
$\sigma$-algebra \emph{generated by $X$} is:
\begin{equation}
    \sigma(X) = \{X^{-1}(B) : B \in \mathcal{S}\} = \{\{X \in B\} : B \in \mathcal{S}\}.
\end{equation}
This is the smallest $\sigma$-algebra on $\Omega$ making $X$ measurable.
\end{definition}

\begin{theorem}[Doob--Dynkin Lemma]
\label{thm:doob_dynkin}
Let $X: \Omega \to S$ and $Y: \Omega \to \R$ be functions. Then $Y$ is
$\sigma(X)$-measurable if and only if $Y = g \circ X$ for some measurable
$g: (S, \mathcal{S}) \to (\R, \Borel(\R))$.
\end{theorem}

\begin{proof}
($\Leftarrow$) If $Y = g(X)$, then $\{Y \leq c\} = \{g(X) \leq c\} = X^{-1}(g^{-1}((-\infty, c])) \in \sigma(X)$.

($\Rightarrow$) We prove this for the case $S = \R^d$ (which covers all RL applications).

\textit{Step 1: Simple functions.} If $Y = \sum a_i \indicator_{A_i}$ with $A_i \in \sigma(X)$, then $A_i = X^{-1}(B_i)$ for some $B_i \in \mathcal{S}$. Define $g = \sum a_i \indicator_{B_i}$. Then $Y = g(X)$.

\textit{Step 2: Non-negative functions.} By Theorem~\ref{thm:approximation}, $Y = \lim_n \varphi_n$ with $\varphi_n$ simple and $\sigma(X)$-measurable. By Step 1, $\varphi_n = g_n(X)$. Define $g = \lim_n g_n$ (well-defined $\mathcal{S}$-a.e.). Then $Y = g(X)$.

\textit{Step 3: General.} Write $Y = Y^+ - Y^-$ and apply Step 2 to each part.
\end{proof}

\begin{rlconnection}[Connection to RL]
The Doob--Dynkin lemma is the formal justification for the pull-out property:
$\E[h(X) \cdot Y \mid X] = h(X) \cdot \E[Y \mid X]$. The function $h(X)$ is
$\sigma(X)$-measurable (it's a function of $X$), so it's ``known'' given $X$
and can be pulled out of the conditional expectation. This is used in every
TD proof when extracting $V(S_t)$ from $\E[\delta_t \mid S_t]$.
\end{rlconnection}

%% ============================================================
%% SECTION 6: CONNECTION TO PROBABILITY
%% ============================================================
\section{From Measures to Probability}
\label{sec:probability}

\subsection{Probability Spaces}
\label{subsec:prob_spaces}

\begin{definition}[Probability Space]
\label{def:prob_space}
A \emph{probability space} is a measure space $(\Omega, \mathcal{F}, \Prob)$ with
$\Prob(\Omega) = 1$. The elements of $\mathcal{F}$ are called \emph{events}, and
$\Prob(A)$ is the \emph{probability} of event $A$.
\end{definition}

\begin{definition}[Random Variable]
\label{def:rv}
A \emph{random variable} on $(\Omega, \mathcal{F}, \Prob)$ is a measurable function
$X: (\Omega, \mathcal{F}) \to (\R, \Borel(\R))$.
\end{definition}

\begin{definition}[Distribution (Pushforward Measure)]
\label{def:distribution}
The \emph{distribution} (or \emph{law}) of a random variable $X$ is the probability
measure $\mu_X$ on $(\R, \Borel(\R))$ defined by:
\begin{equation}
    \mu_X(B) = \Prob(X \in B) = \Prob(X^{-1}(B)) \quad \text{for all } B \in \Borel(\R).
\end{equation}
The \emph{cumulative distribution function} (CDF) is $F_X(x) = \Prob(X \leq x) = \mu_X((-\infty, x])$.
\end{definition}

\begin{proposition}[CDF Determines Distribution]
\label{prop:cdf_determines}
The CDF $F_X$ uniquely determines $\mu_X$.
\end{proposition}

\begin{proof}
The half-lines $\{(-\infty, x] : x \in \R\}$ form a $\pi$-system generating
$\Borel(\R)$. Two measures agreeing on a $\pi$-system agree on the generated
$\sigma$-algebra (Corollary~\ref{cor:uniqueness}). Since $\mu_X((-\infty, x]) = F_X(x)$,
the CDF determines $\mu_X$ on the $\pi$-system, hence on $\Borel(\R)$.
\end{proof}

\subsection{The MDP Probability Space}
\label{subsec:mdp_space}

\begin{example}[Probability Space for a Finite MDP]
\label{ex:mdp_prob_space}
Consider a finite MDP with state space $\mathcal{S}$, action space $\mathcal{A}$,
transition kernel $P$, and policy $\pi$. An infinite trajectory is:
\begin{equation}
    \omega = (s_0, a_0, r_1, s_1, a_1, r_2, s_2, \ldots).
\end{equation}
The probability space is:
\begin{itemize}[nosep]
    \item $\Omega = (\mathcal{S} \times \mathcal{A} \times \R)^\infty$ --- the set of all possible trajectories.
    \item $\mathcal{F} = \sigma(S_0, A_0, R_1, S_1, A_1, R_2, \ldots)$ --- the $\sigma$-algebra generated by all coordinate projections.
    \item $\Prob_\pi$ --- the probability measure determined by the initial distribution,
    policy, and transition kernel via the chain rule:
    \begin{equation}
        \Prob_\pi(S_0 = s_0, A_0 = a_0, S_1 = s_1, \ldots) = p_0(s_0) \prod_{t=0}^{\infty} \pi(a_t | s_t) \, P(s_{t+1} | s_t, a_t).
    \end{equation}
\end{itemize}
The filtration $\mathcal{F}_t = \sigma(S_0, A_0, R_1, \ldots, S_t)$ captures
``information available at time $t$.'' A function $Y$ is $\mathcal{F}_t$-measurable
if and only if $Y = g(S_0, A_0, R_1, \ldots, S_t)$ for some Borel $g$ (Doob--Dynkin).
\end{example}

%% ============================================================
%% SECTION 7: EXERCISES
%% ============================================================
\section{Exercises}
\label{sec:exercises}

\begin{exercise}
Show that $|\mathcal{F}|$ (the number of elements in a finite $\sigma$-algebra on a
finite set $\Omega$) is always a power of 2.
\textit{Hint:} Show that every finite $\sigma$-algebra is generated by a partition of $\Omega$.
\end{exercise}

\begin{exercise}
Prove that $\Borel(\R)$ has the cardinality of the continuum ($|\Borel(\R)| = |\R|$),
while $|\powerset(\R)| = 2^{|\R|} > |\R|$. Conclude that most subsets of $\R$ are not
Borel sets.
\end{exercise}

\begin{exercise}
Let $\mu$ be a finite measure on $(\Omega, \mathcal{F})$ and $A_n \in \mathcal{F}$.
Prove the \emph{first Borel--Cantelli lemma}: if $\sum_n \mu(A_n) < \infty$, then
$\mu(\limsup_n A_n) = 0$, where $\limsup_n A_n = \bigcap_{N=1}^\infty \bigcup_{n=N}^\infty A_n$.
\textit{Hint:} Use continuity from above and the union bound.
\end{exercise}

\begin{exercise}
Let $f, g: (\Omega, \mathcal{F}) \to (\R, \Borel(\R))$ be measurable. Prove that
$\{f < g\}$, $\{f \leq g\}$, $\{f = g\}$ are all in $\mathcal{F}$.
\textit{Hint:} $\{f < g\} = \bigcup_{q \in \Q}(\{f < q\} \cap \{q < g\})$.
\end{exercise}

\begin{exercise}
\label{ex:monotone_class}
Let $\mathcal{A}$ be an algebra on $\Omega$. A \emph{monotone class} is a collection
closed under increasing unions and decreasing intersections. Prove the
\emph{Monotone Class Theorem}: the smallest monotone class containing $\mathcal{A}$
equals $\sigma(\mathcal{A})$.
\textit{Hint:} Adapt the proof of the $\pi$--$\lambda$ theorem.
\end{exercise}

\begin{exercise}
Construct a probability measure on $(\{0,1\}^\N, \mathcal{F})$ corresponding to an
infinite sequence of fair coin flips, where $\mathcal{F}$ is generated by the
cylinder sets. Verify that it satisfies countable additivity.
\end{exercise}

\begin{exercise}[RL Application]
\label{ex:rl_filtration}
Consider an MDP with $\mathcal{S} = \{s_1, s_2, s_3\}$. Write out explicitly all
elements of $\sigma(S_0)$ and $\sigma(S_0, S_1)$. Verify that
$\sigma(S_0) \subseteq \sigma(S_0, S_1)$ (the filtration is increasing).
For a value function $V: \mathcal{S} \to \R$, verify that $V(S_0)$ is
$\sigma(S_0)$-measurable by identifying the function $g$ in the Doob--Dynkin lemma.
\end{exercise}

%% ============================================================
%% SECTION 8: SUMMARY
%% ============================================================
\section{Summary of Key Results}
\label{sec:summary}

\begin{center}
\renewcommand{\arraystretch}{1.4}
\begin{tabular}{p{4cm} p{5.5cm} p{4.5cm}}
\toprule
\textbf{Result} & \textbf{Statement} & \textbf{Role} \\
\midrule
$\sigma$-algebra axioms & Closed under complement and countable union & Framework for events \\
Generated $\sigma$-algebra & $\sigma(\mathcal{C})$ = smallest $\sigma$-algebra $\supseteq \mathcal{C}$ & Defines filtrations \\
$\pi$--$\lambda$ theorem & $\pi$-system $\subseteq$ $\lambda$-system $\implies$ $\sigma(\pi) \subseteq \lambda$ & Uniqueness of measures \\
Measure properties & Monotonicity, sub-additivity, continuity from above/below & All measure arguments \\
Carath\'{e}odory & Pre-measure on algebra extends to $\sigma$-algebra & Constructs Lebesgue measure \\
Generator criterion & Check measurability on generators only & Verifying random variables \\
Approximation theorem & Non-negative measurable $= \lim$ of simple $\uparrow$ & Defines Lebesgue integral \\
Doob--Dynkin & $Y$ is $\sigma(X)$-measurable iff $Y = g(X)$ & Pull-out property \\
\bottomrule
\end{tabular}
\end{center}

\medskip
\noindent\textbf{Next lesson:} Integration theory --- Lebesgue integral, expectation,
monotone/dominated convergence theorems, $L^p$ spaces. These build directly on the
simple function approximation theorem to define $\E[X]$ rigorously.

\end{document}
